\documentclass[12pt]{article}
\usepackage[UTF8]{ctex}
\usepackage{amsmath,amssymb}
\usepackage{braket}
\usepackage{geometry}
\usepackage{fancyhdr}

% 页面设置
\geometry{a4paper, margin=1in}

% 页眉页脚
\pagestyle{fancy}
\fancyhf{}
\fancyhead[L]{量子计算与机器学习第五章作业}
\fancyhead[R]{\today}
\fancyfoot[C]{\thepage}

% 标题信息
\title{\textbf{量子计算与机器学习第五章作业}}
\author{姓名：郝思粤 \quad 学号：PB23151779}
\date{\today}

\begin{document}
\maketitle

\textbf{题目 5.1} 设 $\ket{x}$ 是 $n$ 比特计算基中的一个基矢，
证明
\begin{equation}
H^{\otimes n}\ket{x}
= \frac{1}{\sqrt{2^n}}
\sum_{z\in\{0,1\}^n} (-1)^{x\cdot z} \ket{z},
\end{equation}
其中 $x\cdot z$ 表示按位内积
\(
x\cdot z = \sum_{j=1}^{n} x_j z_j \pmod 2,
\)
$x_j,z_j\in\{0,1\}$ 为 $x,z$ 的第 $j$ 位比特。

\bigskip

\textbf{解：}

首先回顾单比特 Hadamard 门 $H$ 对计算基矢的作用：
\begin{equation}
H\ket{0} = \frac{\ket{0}+\ket{1}}{\sqrt{2}},\qquad
H\ket{1} = \frac{\ket{0}-\ket{1}}{\sqrt{2}}.
\end{equation}
这可以统一写成
\begin{equation}\label{eq:single}
H\ket{x_j} = \frac{\ket{0} + (-1)^{x_j}\ket{1}}{\sqrt{2}},
\qquad x_j\in\{0,1\}.
\end{equation}

令
\(
x = x_1 x_2 \dots x_n
\)
为 $n$ 比特计算基矢对应的比特串，则
\(
\ket{x} = \ket{x_1}\otimes\ket{x_2}\otimes\cdots\otimes\ket{x_n}.
\)
对其施加 $n$ 个 Hadamard 门的张量积 $H^{\otimes n}$，得到
\begin{align}
H^{\otimes n}\ket{x}
&= \bigl(H\ket{x_1}\bigr)\otimes\bigl(H\ket{x_2}\bigr)\otimes\cdots\otimes\bigl(H\ket{x_n}\bigr)\\
&\overset{\eqref{eq:single}}{=}
\left(\frac{\ket{0}+(-1)^{x_1}\ket{1}}{\sqrt{2}}\right)
\otimes
\left(\frac{\ket{0}+(-1)^{x_2}\ket{1}}{\sqrt{2}}\right)
\otimes \cdots \otimes
\left(\frac{\ket{0}+(-1)^{x_n}\ket{1}}{\sqrt{2}}\right)\\
&=
\frac{1}{\sqrt{2^n}}
\left(\ket{0}+(-1)^{x_1}\ket{1}\right)
\otimes
\left(\ket{0}+(-1)^{x_2}\ket{1}\right)
\otimes\cdots\otimes
\left(\ket{0}+(-1)^{x_n}\ket{1}\right).
\end{align}

把上式按张量积逐项展开。  
每一项都对应一个比特串
\(
z=z_1 z_2\dots z_n\in\{0,1\}^n,
\)
其中
\(
z_j=0
\)
表示在第 $j$ 个括号中选择 $\ket{0}$，
而
\(
z_j=1
\)
表示选择 $(-1)^{x_j}\ket{1}$。  
因此，展开后的一般项为
\begin{equation}
\prod_{j=1}^{n}(-1)^{x_j z_j} \ket{z_1}\otimes\ket{z_2}\otimes\cdots\otimes\ket{z_n}
= (-1)^{\sum_{j=1}^{n} x_j z_j}\ket{z},
\end{equation}
其中
\(
\ket{z}=\ket{z_1 z_2\dots z_n}.
\)

注意指数
\(
\sum_{j=1}^{n} x_j z_j
\)
只在模 $2$ 意义下取值（因为 $x_j,z_j\in\{0,1\}$ 且 $(-1)^k$ 只依赖 $k$ 的奇偶性），
这恰好就是题目中的按位内积 $x\cdot z$。  
于是整个展开可写作
\begin{equation}
H^{\otimes n}\ket{x}
=
\frac{1}{\sqrt{2^n}}
\sum_{z\in\{0,1\}^n}
(-1)^{x\cdot z} \ket{z},
\end{equation}
这正是要证明的式子。

\hfill$\blacksquare$

\end{document}
